
\section{有限群作用与商空间}
\label{sec:a1-3}

下面定义有限群在带仓西结构空间上的作用\liftIndex{I318}及其商空间。

\begin{definition}{A1.45}\label{def:a1-45}
设 $\varphi\colon X\to X$ 是带仓西结构空间 $X$ 的同胚。若下列条件
成立，就称它诱导一个\emph{仓西结构自同构}。取 $p\in X$，令
$p'=\varphi(p)$。设
\[
 (V_p,E_p,\Gamma_p,\psi_p,s_p),
 \qquad
 (V_{p'},E_{p'},\Gamma_{p'},\psi_{p'},s_{p'})
\]
分别是 $p$ 与 $p'$ 的仓西邻域。要求存在
\[
 \rho_p\colon\Gamma_p\longrightarrow\Gamma_{p'},
 \qquad
 \varphi_p\colon V_p\longrightarrow V_{p'},
 \qquad
 \widehat\varphi_p\colon E_p\times V_p
       \longrightarrow E_{p'}\times V_{p'},
\]
使得：
\begin{description}[style=multiline,leftmargin=7.2em,labelwidth=6.5em]
\item[\textup{(A1.46.1)}] $\rho_p$ 是群同构；
\item[\textup{(A1.46.2)}] $\varphi_p$ 是 $\rho_p$-等变微分同胚；
\item[\textup{(A1.46.3)}] $\widehat\varphi_p$ 是覆盖 $\varphi_p$ 的
$\rho_p$-等变向量丛同构；
\item[\textup{(A1.46.4)}]
$s_{p'}\circ\varphi_p=\widehat\varphi_p\circ s_p$；
\item[\textup{(A1.46.5)}] $\varphi_p$ 在 $s_p^{-1}(0)$ 上的限制诱导同胚
\[
 s_p^{-1}(0)/\Gamma_p\longrightarrow
 s_{p'}^{-1}(0)/\Gamma_{p'},
\]
把它记作 $\underline\varphi_p$。此时
$\psi_{p'}\circ\underline\varphi_p=\varphi\circ\psi_p$。
\end{description}

还要求 $\rho_p,\varphi_p,\widehat\varphi_p$ 与仓西结构的坐标变换相容，
意义如下。设
\[
 q\in\psi_p(s_p^{-1}(0)/\Gamma_p),
 \qquad
 q'\in\psi_{p'}(s_{p'}^{-1}(0)/\Gamma_{p'}),
 \qquad \varphi(q)=q'.
\]
令
\[
 (\widehat\phi_{pq},\phi_{pq},h_{pq}),
 \qquad
 (\widehat\phi_{p'q'},\phi_{p'q'},h_{p'q'})
\]
为相应的坐标变换。要求存在
$\gamma_{pqp'q'}\in\Gamma_{p'}$，使得：
\begin{description}[style=multiline,leftmargin=7.2em,labelwidth=6.5em]
\item[\textup{(A1.47.1)}]
\[
 \rho_p\circ h_{pq}
 =\gamma_{pqp'q'}\cdot(h_{p'q'}\circ\rho_q)\cdot
  \gamma_{pqp'q'}^{-1};
\]
\item[\textup{(A1.47.2)}]
\[
 \varphi_p\circ\phi_{pq}
 =\gamma_{pqp'q'}\cdot(\phi_{p'q'}\circ\varphi_q);
\]
\item[\textup{(A1.47.3)}]
\[
 \widehat\varphi_p\circ\widehat\phi_{pq}
 =\gamma_{pqp'q'}\cdot
  (\widehat\phi_{p'q'}\circ\widehat\varphi_q).
\]
\end{description}
我们把 $((\rho_p,\varphi_p,\widehat\varphi_p)_p;\varphi)$ 称为这一
仓西结构的自同构。

若两个自同构
\[
 ((\rho_p,\varphi_p,\widehat\varphi_p)_p;\varphi),
 \qquad
 ((\rho'_p,\varphi'_p,\widehat\varphi'_p)_p;\varphi')
\]
满足 $\varphi=\varphi'$，并且对每个 $p$ 都存在
$\gamma_p\in\Gamma_{\varphi(p)}$，使得
\begin{description}[style=multiline,leftmargin=7.2em,labelwidth=6.5em]
\item[\textup{(A1.48.1)}]
$\rho'_p=\gamma_p\cdot\rho_p\cdot\gamma_p^{-1}$；
\item[\textup{(A1.48.2)}]
$\varphi'_p=\gamma_p\cdot\varphi_p$；
\item[\textup{(A1.48.3)}]
$\widehat\varphi'_p=\gamma_p\cdot\widehat\varphi_p$，
\end{description}
就称这两个自同构\emph{共轭}\translatornote{源书最初把 $\widehat\varphi_p$ 的型别仅写成
$E_p\to E_{p'}$，但条件 (A1.46.3) 要求它覆盖底映射 $\varphi_p$；依
定义 A1.1 的平凡障碍丛及引理 A1.49 的同型提升，补全为
$E_p\times V_p\to E_{p'}\times V_{p'}$。}。
\end{definition}

两个自同构的复合定义为
\begin{align*}
 &((\rho_p^1,\varphi_p^1,\widehat\varphi_p^1)_p;\varphi^1)
 \circ
 ((\rho_p^2,\varphi_p^2,\widehat\varphi_p^2)_p;\varphi^2)
 \\
 &\quad=
 \Bigl((
 \rho_{\varphi^2(p)}^1\circ\rho_p^2,
 \varphi_{\varphi^2(p)}^1\circ\varphi_p^2,
 \widehat\varphi_{\varphi^2(p)}^1\circ\widehat\varphi_p^2
 )_p;
 \varphi^1\circ\varphi^2\Bigr).
\end{align*}
容易验证：若左端按 (A1.47) 的意义与坐标变换相容，则右端也相容；复合
也与上述共轭关系相容。

若该提升与
\[
 \Lambda^{\mathrm{top}}E_p^*\otimes\Lambda^{\mathrm{top}}TV_p
\]
的平凡化相容，就称它\emph{保定向}。以 $\operatorname{Aut}(X)$ 表示
$X$ 的全部自同构之共轭类所成集合，以 $\operatorname{Aut}_0(X)$ 表示
全部保定向自同构之共轭类所成集合。复合使二者成为群。

设有限群 $G$ 作用在紧空间 $X$ 上，并假定 $X$ 有仓西结构。若对每个
$g\in G$，同胚 $x\mapsto gx$ 都提升为仓西结构的一个自同构 $g_*$，
并且 $g_*$ 与 $h_*$ 的复合共轭于 $(gh)_*$，就称 $G$ 作用在 $X$ 上
（把 $X$ 看作带仓西结构空间）。换言之，$G$ 在 $X$ 上的作用就是群同态
\[
 G\longrightarrow\operatorname{Aut}(X).
\]
带仓西结构空间上的一个\emph{对合}就是一个 $\mathbb Z_2$-作用。

\begin{lemma}{A1.49}\label{lem:a1-49}
若有限群 $G$ 作用在带仓西结构空间 $X$ 上，则商空间 $X/G$ 有仓西结构。

若 $X$ 有切丛且该作用保持切丛，则商空间有切丛。若 $X$ 已定向且该作用
保持定向，则商空间已定向。
\end{lemma}

\begin{proof}
先假设 $G$ 的作用有效。取 $g\in G$、$p\in X$，并设
$(V_p,E_p,\Gamma_p,\psi_p,s_p)$ 是 $p$ 的仓西邻域。置
\[
 G_p=\{g\in G\mid gp=p\}.
\]
依定义，$g\in G_p$ 诱导
\[
 \varphi_{g,p}\colon V_p\longrightarrow V_p,
 \qquad
 \widehat\varphi_{g,p}\colon E_p\times V_p
     \longrightarrow E_p\times V_p,
 \qquad
 \rho_{g,p}\colon\Gamma_p\longrightarrow\Gamma_p.
\]
把 $G_p$ 与 $\{\varphi_{g,p}\mid g\in G_p\}$ 识别。

令 $\Gamma_{[p]}$ 是 $V_p$ 的微分同胚群中由 $G_p$ 与 $\Gamma_p$
生成的群。我们断言，$\Gamma_{[p]}$ 的每个元素都唯一地写成复合
$g\circ\gamma$，其中 $g\in G_p$、$\gamma\in\Gamma_p$。证明如下。
对 $g_1,g_2\in G_p$，取 $\gamma_{g_1,g_2}\in\Gamma_p$，使
\[
 g_1(g_2(x))=\gamma_{g_1,g_2}((g_1g_2)(x)).
\]
这样的 $\gamma_{g_1,g_2}$ 存在，因为 $g_1$ 与 $g_2$ 的复合共轭于
$g_1g_2$。于是 $\Gamma_{[p]}$ 上的乘法 $\circ$ 由
\[
 g_1\circ g_2=\gamma_{g_1,g_2}\circ(g_1g_2),
 \qquad
 g\circ\gamma=\rho_{g,p}(\gamma)\circ g
\]
给出，其中 $g,g_1,g_2\in G_p$、$\gamma\in\Gamma_p$。这就证明了断言。

由此，$\Gamma_{[p]}$ 是有限群。此外，有正合列
\begin{equation}
 1\longrightarrow\Gamma_p\longrightarrow\Gamma_{[p]}
 \longrightarrow G_p\longrightarrow1.
\tag{A1.50}\label{eq:a1-50}
\end{equation}
$G_p$ 的每个元素由式~\eqref{eq:a1-50} 在 $\Gamma_p$ 上诱导的作用，
由 $g\mapsto\rho_{g,p}$ 给出。（注意，$g\mapsto\rho_{g,p}$ 不是群同态。）
正合列~\eqref{eq:a1-50} 未必分裂。事实上，$\gamma_{g_1,g_2}$ 对应于
决定群扩张~\eqref{eq:a1-50} 的 $2$-闭上链。

可以把 $\Gamma_{[p]}$ 在 $V_p$ 上的作用提升为它在
$E_p\times V_p$ 上的作用，因为 $\widehat\varphi_{g,p}$（$g\in G_p$）
与 $\gamma\in\Gamma_p$ 之间的关系由同样的 $\rho_{g,p}$ 和
$\gamma_{g_1,g_2}$ 描述。仓西映射 $s_p\colon V_p\to E_p$ 是
$\Gamma_{[p]}$-等变的。由于
\[
 V_p/\Gamma_{[p]}\cong(V_p/\Gamma_p)/G_p,
\]
映射 $\psi_p\colon s_p^{-1}(0)/\Gamma_p\to X$ 诱导
\[
 \psi_{[p]}\colon s_p^{-1}(0)/\Gamma_{[p]}\longrightarrow X/G,
\]
它是到 $[p]$ 的一个邻域上的同胚。因此，我们构造出了
$[p]\in X/G$ 的仓西邻域
\[
 (V_p,E_p,\Gamma_{[p]},\psi_{[p]},s_p).
\]
直接构造这些仓西邻域之间的坐标变换，并验证定义~\ref{def:a1-3} 与
定义~\ref{def:a1-5} 所要求的性质。关于切丛与定向的结论，由它们的
定义立即得到。

若 $G$ 的作用无效，就固定一个有效的有限维表示
\[
 G\longrightarrow GL(E).
\]
把仓西邻域 $(V_p,E_p,\Gamma_p,\psi_p,s_p)$ 替换为
\[
 (V_p\times E,E_p\times E,\Gamma_p,\psi_p,s_p\oplus\operatorname{id}).
\]
此时 $G$ 在这一仓西邻域上的作用有效，因而可应用上述构造。引理
\ref{lem:a1-49} 得证\translatornote{源书在提升 $\Gamma_{[p]}$ 的作用时把已定义的
$\rho_{g,p}$ 写成 $\rho_{p,g}$；译文依前一页的定义、乘法公式及
式 (A1.50) 后的同一记号恢复。}。
\end{proof}

\begin{remark}{A1.51}\label{rem:a1-51}
\textup{(1)} 在例~\ref{ex:a1-7} 的情形中，假设存在一个群 $G'$，它以
$G$ 为有限指数正规子群；又假设 $G$ 在 $W_1,W_2$ 上的作用可延拓为
$G'$-作用，并且 $F$ 是 $G'$-等变的。则可选取
$X=F^{-1}(0)/G$ 上的仓西结构，使它成为 $H$-等变的，其中
$H=G'/G$。此外，引理~\ref{lem:a1-49} 给出的 $X/H$ 上的商仓西结构，
同 $F^{-1}(0)/G'$ 上的仓西结构同构。

\textup{(2)} 群作用的定义可显然推广到连续群；见
\textbf{[FuOn99II]} 第~22 节。此时，在作用恰当且稳定子群有限的假设下，
引理~\ref{lem:a1-49} 仍可推广。本书不需要这一推广，所以不作证明。
\end{remark}
